\documentclass[11pt,fontset=fandol]{ctexart}
\usepackage[a4paper,margin=1in]{geometry}
\usepackage{amsmath,amssymb,amsfonts,mathtools}
\usepackage{booktabs}
\usepackage{graphicx}
\usepackage{caption}
\usepackage{enumitem}
\usepackage{array}
\usepackage{float}
\usepackage[hidelinks]{hyperref}
\usepackage{xurl}
\usepackage{microtype}
\usepackage{longtable}

\title{Exact Recursion 方法教学主稿\\\large 从最基础的随机游走公式，到主线 one-target 机制稿中的实际计算}
\author{}
\date{\today}

\setlist{nosep}
\setlength{\textfloatsep}{8pt plus 2pt minus 2pt}
\setlength{\floatsep}{6pt plus 2pt minus 2pt}
\setlength{\intextsep}{6pt plus 2pt minus 2pt}
\setlength{\abovecaptionskip}{4pt}
\setlength{\belowcaptionskip}{0pt}
\setlength{\tabcolsep}{5pt}
\renewcommand{\arraystretch}{1.12}
\setlength{\emergencystretch}{2em}

\newcommand{\GatingFigDir}{../../grid2d_one_two_target_gating/artifacts/figures}
\newcommand{\BenchFigDir}{../../luca_vs_recursion_unified_benchmark/artifacts/figures}
\newcommand{\RingFigDir}{../../ring_lazy_flux/artifacts/figures}

\begin{document}
\maketitle

\section*{结论先行}
这份文稿想做一件很具体的事：把本仓库主线中真正大量使用的 \emph{Exact Recursion}（也常叫 time recursion、时域递推、精确吸收递推）从头讲清楚，而且假设读者只记得非常基础的线性代数和微积分。

\paragraph{先说最短答案。}
在你的研究里，最重要的不是“随机挑很多 Monte Carlo 轨迹看图像”，也不是“默认先把所有问题都丢进 Giuggioli/Sarvaharman 一类的生成函数 family 里求解”。真正的主求解器，尤其是在当前二维主线里，是：先把问题写成一个有限状态的吸收 Markov 链，再把\textbf{尚未命中目标的概率质量}一步一步精确推进。每一步流入目标点的那部分概率，就是这一时刻的首达时间概率。

\paragraph{为什么这个方法重要。}
它重要有三个原因：
\begin{enumerate}
  \item 它直接输出完整的首达时间分布 $f(t)$，所以你能看到早峰、谷底和晚峰；
  \item 它不是 Monte Carlo 采样，因此没有统计噪声；只要转移概率写对了，给定时间窗内的结果就是确定性的；
  \item 它和你的机制分析天然兼容，因为同一套递推不仅能给总曲线，还能给通道分解、窗口条件化、timing observables 和 budget observables。
\end{enumerate}

\paragraph{这份稿子和仓库主线的关系。}
本文会一直拿 \path{grid2d_one_two_target_gating} 的 one-target shared symmetric baseline 当主例子。因为在当前仓库语言里，最核心的一句话已经很明确：标准配置是 shared symmetric baseline；\texttt{peak1} 是 corridor fast branch；\texttt{peak2} 的主体是 left-open detour/rollback branch；膜穿越只是慢支路的次级修饰，而不是主角。换句话说，Exact Recursion 在这里不只是“算了一条曲线”，而是把整套机制语言撑起来了。

\begin{figure}[H]
  \centering
  \includegraphics[width=0.88\linewidth]{\GatingFigDir/one_target_gate_geometry.pdf}
  \caption{当前主线 one-target 例子的标准几何图。本文会一直围绕这张图来解释：状态空间是什么、目标是什么、哪些区域对应快支路、哪些区域对应慢支路。}
  \label{fig:main-geometry}
\end{figure}

\section{最基础的起点：什么是离散时间随机游走}
\paragraph{先把问题说得最朴素一点。}
我们关心一个粒子（或者 walker）在离散的时间步 $t=0,1,2,\dots$ 上运动。每一步，它会以某些概率走到相邻格点，或者在 lazy 模型里原地停留。把“当前位置”记作 $X_t$，这就是一个最基本的离散时间随机过程。

\paragraph{什么是状态。}
对二维格点模型来说，一个状态通常就是一个格点坐标，例如 $(x,y)$。如果格子大小是 $N\times N$，那全部状态就是所有合法的 $(x,y)$。如果后面我们再给轨迹附加标签，例如“是否已经 leak”“是否已经回到起始盆地”，那状态空间会更大，但思想完全一样：\textbf{状态就是系统此刻所处的全部信息。}

\paragraph{什么是转移概率。}
如果当前在状态 $i$，下一步走到状态 $j$ 的概率记作 $P_{ij}$。所有这些概率拼起来就是转移矩阵 $P$。这里最容易忘掉的一点是：矩阵里的每个元素不是一个数列，而是“一步转移”的概率。也就是说，矩阵负责回答“下一步怎么走”，而不是“最终会不会命中”。

\paragraph{什么是 Markov 性。}
Markov 性的意思是：在知道当前状态以后，下一步怎么走只依赖当前状态，不依赖更早以前的历史。它并不意味着系统“没有记忆”，而是意味着只要你把需要的记忆都塞进状态定义里，\textbf{当前状态就已经包含了未来一步所需的全部信息。}

\paragraph{二维随机游走的最基本主方程。}
先看最简单的、无偏、非 lazy 的二维随机游走。若
\[
  p_t(x,y)=\mathbb P(X_t=(x,y)),
\]
那么在内部点上有
\begin{equation}
  p_{t+1}(x,y)
  =
  \frac14 p_t(x-1,y)
  +
  \frac14 p_t(x+1,y)
  +
  \frac14 p_t(x,y-1)
  +
  \frac14 p_t(x,y+1).
  \label{eq:basic-master}
\end{equation}
这条式子读法很简单：要想在第 $t+1$ 步站到 $(x,y)$，那第 $t$ 步必须先在四个邻居之一，而且从那个邻居走到这里。四条可能路径的概率加起来，就是新的位置分布。

\paragraph{为什么矩阵乘法会自然出现。}
如果把所有格点按某个固定顺序排成一个长向量
\[
  p_t = \bigl(p_t(s_1),p_t(s_2),\dots,p_t(s_M)\bigr),
\]
那上面的每个“从哪些邻居流进来”的式子其实都可以统一写成
\begin{equation}
  p_{t+1}=p_t P
  \quad\text{或}\quad
  p_{t+1}=P^\top p_t,
  \label{eq:matrix-evolution}
\end{equation}
具体写成左乘还是右乘，只取决于你把 $P_{ij}$ 定义成“从 $i$ 到 $j$”还是“从 $j$ 到 $i$”。仓库里两种写法都出现过，但物理内容是一样的：\textbf{矩阵乘法只是把很多个“流入求和”方程压缩成一条紧凑写法。}

\paragraph{lazy 和 bias 在这里怎么进来。}
如果模型是 lazy 的，那每一步除了向邻点走，还允许以一定概率原地停留。最常见的写法是引入参数 $q\in(0,1]$ 表示“本步真的去移动的总概率”，那么原地停留概率就是 $1-q$。如果再引入方向偏置，例如向右更容易，那左右上下四个方向的权重就不再均匀，而是按照全局偏置、局部箭头、sticky site 和 barrier/membrane 等规则重新分配。仓库里的二维主线就是这种“lazy + bias + 局部结构”的一般形式。

\section{从位置分布到首达时间}
\paragraph{位置分布不是我们真正想要的量。}
只知道“第 $t$ 步时粒子在什么地方”的分布还不够。你的研究真正关心的是：\textbf{粒子第一次命中目标是在什么时候。} 这就是首达时间（first-passage time）。

\paragraph{首达时间怎么定义。}
设目标集合为 $T_{\mathrm{set}}$。首达时间定义为
\begin{equation}
  T
  =
  \inf\{t\ge 1: X_t\in T_{\mathrm{set}}\}.
  \label{eq:fpt-def}
\end{equation}
如果目标只有一个格点 $(x_\star,y_\star)$，那就是“第一次走到这个目标格点的时刻”。如果目标是一条对角线、两个 target 的并集，或者 pair chain 里的 encounter set，定义形式完全一样，只是 $T_{\mathrm{set}}$ 更大。

\paragraph{三个最重要的分布量。}
一旦定义了 $T$，我们就能写出：
\begin{align}
  f(t) &= \mathbb P(T=t), \label{eq:fpt-pmf}\\
  F(t) &= \mathbb P(T\le t), \label{eq:fpt-cdf}\\
  S(t) &= \mathbb P(T>t). \label{eq:survival-def}
\end{align}
其中 $f(t)$ 是首达时间的 pmf，也就是你论文里真正画双峰的那条曲线；$F(t)$ 是累计命中概率；$S(t)$ 是到第 $t$ 步为止还没命中的概率，也叫 survival。

\paragraph{为什么“在目标点上的概率”不等于“首次命中的概率”。}
这是整篇文稿最关键的概念之一，值得重复说。

若你只是看普通位置分布 $p_t(\text{target})$，那它统计的是“第 $t$ 步在目标上”的概率。这会把两种轨迹混在一起：
\begin{itemize}
  \item 恰好在第 $t$ 步第一次命中目标的轨迹；
  \item 更早就已经到过目标、后来还继续算在目标上的轨迹。
\end{itemize}
而首达概率 $f(t)=\mathbb P(T=t)$ 只允许第一类轨迹，绝不允许第二类轨迹重复计数。所以 first-passage 的核心工作，不是普通传播，而是要把已经命中过的概率质量\textbf{从后续传播里剔除出去。}

\paragraph{吸收态和瞬态态。}
最自然的做法是把目标当成吸收态。所谓吸收态，就是一旦到达就不再离开。与之相对，其余还会继续传播的状态叫瞬态态（transient states）。从 first-passage 的角度看，我们真正要推进的是瞬态态上的概率质量，因为一旦某条路径第一次进入吸收态，它的首达时间已经定下来了，不应该再参与后续的“未命中传播”。

\section{Exact Recursion 的核心想法}
\paragraph{这套方法到底推进什么。}
Exact Recursion 推进的不是普通位置分布，而是
\begin{equation}
  u_t(s)=\mathbb P(X_t=s,\;T>t),
  \label{eq:unabsorbed-dist}
\end{equation}
也就是：“到第 $t$ 步为止还没有命中目标，并且此刻位于状态 $s$ 的概率。”

\paragraph{为什么这个量更合适。}
因为一旦我们知道了所有状态上的未吸收质量 $u_t(s)$，就能同时拿到三件东西：
\begin{enumerate}
  \item 下一步还没命中的分布 $u_{t+1}$；
  \item 下一步首次命中的概率 $f(t+1)$；
  \item 当前 survival $S(t)=\sum_s u_t(s)$。
\end{enumerate}
也就是说，\textbf{同一份递推同时控制了曲线和生存函数。}

\paragraph{把状态空间拆成两部分。}
把状态排成“瞬态态在前，吸收态在后”的顺序后，转移矩阵可以分块写成
\begin{equation}
  P=
  \begin{pmatrix}
    Q & R \\
    0 & I
  \end{pmatrix}.
  \label{eq:block-matrix}
\end{equation}
这里：
\begin{itemize}
  \item $Q$ 是瞬态态到瞬态态的一步转移；
  \item $R$ 是瞬态态到吸收态的一步流量；
  \item 右下角的 $I$ 表示一旦进入吸收态就留在那里。
\end{itemize}
对单目标问题，$R$ 其实可以压成一列向量 $r$，因为“吸收态”只有一个目标点。

\paragraph{最核心的递推公式。}
若当前未吸收分布写成行向量 $u_t$，那么 Exact Recursion 的母式就是
\begin{align}
  u_{t+1} &= u_t Q, \label{eq:exact-main-u}\\
  f(t+1) &= u_t r, \label{eq:exact-main-f}\\
  S(t) &= u_t\mathbf 1. \label{eq:exact-main-s}
\end{align}
这三条式子就是整套方法的真正核心。后面你在二维主线里看到的一切：总曲线、双峰窗口、通道质量、window-conditioned observables，都是建立在它之上的。

\paragraph{每条式子在说什么。}
\begin{itemize}
  \item 式 \eqref{eq:exact-main-u} 说：下一步仍然未命中的概率质量，只来自当前未命中的质量通过瞬态到瞬态的边继续流动。
  \item 式 \eqref{eq:exact-main-f} 说：下一步首次命中的质量，就是当前未命中质量通过瞬态到目标的边流出去的总量。
  \item 式 \eqref{eq:exact-main-s} 说：survival 就是当前所有未命中质量加总。
\end{itemize}

\paragraph{为什么这三条式子是严格正确的。}
原因其实只用到两件事：全概率公式和 Markov 性。

对任一瞬态态 $j$，有
\[
  u_{t+1}(j)
  =
  \mathbb P(X_{t+1}=j,\;T>t+1).
\]
要想在第 $t+1$ 步还没命中而且站在 $j$，那么第 $t$ 步一定已经处于某个瞬态态 $i$，并且当时还没命中。由 Markov 性，下一步从 $i$ 走到 $j$ 的概率只取决于 $i$ 本身，所以
\[
  u_{t+1}(j)=\sum_i \mathbb P(X_t=i,\;T>t)\,Q_{ij}
  =
  \sum_i u_t(i)Q_{ij}.
\]
这就是 \eqref{eq:exact-main-u}。

同理，若 $r_i$ 表示“从瞬态态 $i$ 下一步直接命中目标”的概率，那么
\[
  f(t+1)=\sum_i u_t(i)r_i,
\]
这就是 \eqref{eq:exact-main-f}。所以它不是近似方法，而是在\textbf{精确统计所有长度为 $t$ 的未命中轨迹，在下一步流向哪里。}

\paragraph{把它写成矩阵幂。}
从 \eqref{eq:exact-main-u} 出发反复递推可得
\begin{equation}
  u_t = u_0 Q^t.
  \label{eq:ut-power}
\end{equation}
因此
\begin{equation}
  f(t+1)=u_0Q^t r.
  \label{eq:f-power}
\end{equation}
如果你更习惯列向量写法，也可以写成
\[
  p_{t+1}=Qp_t,\qquad f(t)=r^\top p_t,
\]
这就是仓库里多个文稿使用的等价口径。

\subsection*{为什么“每步把 target 概率拿掉”正是 first-passage 的关键}
这句话很值得单独成节，因为你最近反复问到它。

\paragraph{先想一个错误做法。}
如果我们每一步都按普通位置分布推进，但不把目标点概率拿掉，那么一条在第 $5$ 步第一次命中目标的轨迹，在第 $6,7,8,\dots$ 步还会继续留在目标点上。这样一来，后续每一步的“在目标上概率”都会把它反复算进去。

\paragraph{这为什么错。}
因为首达时间只允许每条轨迹在\textbf{第一次}命中时记账一次。之后它已经完成任务了，不能再重复贡献给 $f(t)$。

\paragraph{正确做法是什么。}
正确做法就是：
\begin{enumerate}
  \item 先让当前未命中质量按一步转移走出去；
  \item 统计其中有多少流进了目标，这一部分就是当前的首达质量；
  \item 把流进目标的这部分从后续传播里拿掉；
  \item 剩下的才是新的未吸收质量，继续下一轮递推。
\end{enumerate}
所以“拿掉 target 概率”不是一个额外技巧，而是\textbf{把普通传播改造成 first-passage 传播的本质步骤。}

\paragraph{为什么这仍然是精确的。}
因为你并没有扔掉任何有效信息。你只是把“已经完成首达的质量”从“尚未首达的传播池”里移出，并把它记到正确的时间桶里。对于 first-passage 问题，这恰恰是\textbf{唯一正确的记账方式。}

\subsection*{为什么 Exact Recursion 不是 Monte Carlo}
MC 的做法是抽样很多条随机轨迹，再用频率近似概率。样本数有限时，曲线会有统计噪声，尤其尾部更明显。Exact Recursion 完全不同：它不是抽样，而是直接对概率质量做确定性推进。你可以把它看成“把指数多条轨迹按当前落点无损压缩并求和”。只要状态空间是有限的、转移概率是已知的，那么在给定时间窗内它给出的 $f(t)$ 就是精确的，误差只来自浮点和你主动设定的时间截断。

\section{一个极小 toy example}
\paragraph{为什么一定要先看小例子。}
上面的公式已经是完整正确的，但如果第一次接触 first-passage，最容易卡住的地方不是符号，而是“概率质量到底在怎么流”。所以我们先看一个非常小的图模型，只保留三个状态：
\[
  A \longleftrightarrow B \longrightarrow T,
\]
其中 $T$ 是目标吸收态，$A,B$ 是瞬态态。

\paragraph{具体转移规则。}
设从 $A$ 出发：
\[
  A\to A:0.2,\qquad A\to B:0.8,\qquad A\to T:0.
\]
从 $B$ 出发：
\[
  B\to A:0.3,\qquad B\to B:0.1,\qquad B\to T:0.6.
\]
若初始位置在 $A$，那么
\[
  u_0=(1,0).
\]
瞬态子矩阵和吸收流量向量分别是
\[
  Q=
  \begin{pmatrix}
    0.2 & 0.8\\
    0.3 & 0.1
  \end{pmatrix},
  \qquad
  r=
  \begin{pmatrix}
    0\\
    0.6
  \end{pmatrix}.
\]

\paragraph{第 1 步。}
由 $u_1=u_0Q$ 得
\[
  u_1=(1,0)
  \begin{pmatrix}
    0.2 & 0.8\\
    0.3 & 0.1
  \end{pmatrix}
  =(0.2,0.8).
\]
这表示：走了 1 步以后，仍然未命中的质量里，有 $0.2$ 在 $A$，有 $0.8$ 在 $B$。

而首达概率
\[
  f(1)=u_0r=(1,0)
  \begin{pmatrix}
    0\\
    0.6
  \end{pmatrix}
  =0.
\]
因为从 $A$ 不能一步直接到目标。

\paragraph{第 2 步。}
再递推一次：
\[
  u_2=u_1Q=(0.2,0.8)
  \begin{pmatrix}
    0.2 & 0.8\\
    0.3 & 0.1
  \end{pmatrix}
  =(0.28,0.24).
\]
这意味着到第 2 步以后，仍未命中的质量剩下 $0.52$。

这一轮的首达概率是
\[
  f(2)=u_1r=(0.2,0.8)
  \begin{pmatrix}
    0\\
    0.6
  \end{pmatrix}
  =0.48.
\]
也就是说，有 $48\%$ 的路径是在第 2 步第一次命中目标。

\paragraph{survival 怎么看。}
第 1 步 survival 是
\[
  S(1)=0.2+0.8=1.
\]
第 2 步 survival 是
\[
  S(2)=0.28+0.24=0.52.
\]
和首达质量的关系也很直观：
\[
  S(2)=S(1)-f(2)=1-0.48=0.52.
\]
这正是“流进目标的质量被拿掉以后，剩余的未命中质量继续传播”。

\paragraph{这个小例子告诉了我们什么。}
它告诉我们：Exact Recursion 并不是在“追踪某一条轨迹”，而是在追踪\textbf{所有未命中轨迹的总概率质量怎样在瞬态态空间里流动。} 只要你理解了这个 toy example，再回头看二维格点，其实只是状态数更多、边更多、几何更复杂，但数学结构没有变。

\section{映射到仓库代码：它到底是怎么实现的}
\paragraph{第一步：先把几何转成一步转移。}
仓库不会直接在 LaTeX 里求概率，而是先把几何、bias、sticky、barrier、membrane 等规则全部转成一步转移表。对二维主线，核心文件是
\begin{quote}
\path{packages/vkcore/src/vkcore/grid2d/model_core_reflecting.py}
\end{quote}
它做的第一件事，是对每个格点列出所有允许的一步跳转，并附上对应概率。这样一来，“物理几何”就被翻译成了“离散状态图 + 边权重”。

\paragraph{第二步：把目标边拆出去。}
同一个文件里，\path{build_exact_arrays} 会把所有非目标格点收集成瞬态态列表，然后对每个瞬态态 $i$ 做两类拆分：
\begin{itemize}
  \item 如果一步会走到目标，就把这条边的概率累积到 $r_i$；
  \item 如果一步还是走到非目标格点，就把这条边记到 \path{src\_idx/dst\_idx/probs} 里。
\end{itemize}
这一步正是在把整个一步转移矩阵分解成“瞬态到瞬态的 $Q$”和“瞬态到目标的 $r$”。

\paragraph{第三步：做稀疏推进。}
在 one-target 的通用实现里，主循环大致等价于：
\begin{quote}
当前未命中分布 $p$ 先按照所有瞬态边推进到 $p_{\text{next}}$；\\
读出其中流入目标的质量作为 \path{hit\_total}；\\
把目标位置清零；\\
剩余部分作为下一轮继续递推的未吸收质量。
\end{quote}
仓库中的通用 one-target 实现可以在
\begin{quote}
\path{packages/vkcore/src/vkcore/grid2d/rect_bimodality/cli.py}
\end{quote}
里的 \path{run_exact_one_target_rect} 找到。

\paragraph{为什么这等价于稀疏矩阵乘法。}
如果真的把 $Q$ 显式存成一个大矩阵，每一轮只是在做一次向量乘矩阵：
\[
  u_{t+1}=u_tQ.
\]
代码里不显式存整块 $Q$，而是只存非零边：
\[
  \texttt{src\_idx},\ \texttt{dst\_idx},\ \texttt{probs}.
\]
然后用“把每条边贡献加到目标节点”的方式完成同样的求和。这其实就是稀疏矩阵乘法的散点累加写法。它的好处是：状态多的时候不必浪费空间存大量零元。

\paragraph{为什么说这仍然是 exact。}
这里的“稀疏”只是存储形式稀疏，不是数学近似稀疏。只要所有非零边都完整地被枚举出来，那么这和显式乘上完整矩阵 $Q$ 的结果完全一样。换句话说，\textbf{稀疏实现改变的是速度和内存，不改变数学答案。}

\section{主线真实例子：shared symmetric one-target baseline}
\paragraph{为什么选这个例子。}
因为它是你当前仓库 one-target 主语言的中心。它不是某个极端 showcase，而是标准配置。主线报告里已经明确说了：canonical one-target 语言回到了 shared symmetric baseline，第一峰是 corridor fast branch，第二峰的主体是 left-open detour/rollback branch，膜穿越只是慢支路的次级修饰。

\paragraph{这个例子里状态空间长什么样。}
主图见图 \ref{fig:main-geometry}。如果用最朴素的话来说，这个例子有五层结构：
\begin{enumerate}
  \item 一个从左向右的 corridor，是最快的主通道；
  \item corridor 上下都有 outside 区域，可以形成更长的绕行；
  \item corridor 与 outside 的交界上有半透膜，它会改变进出 outside 的难易；
  \item 左端附近有一个起始盆地 $A$，轨迹是否离开又返回 $A$ 会决定 rollback 类别；
  \item 右端有 target，是唯一吸收点。
\end{enumerate}

\paragraph{Exact Recursion 在这里具体输出什么。}
最基本的输出是首达时间 pmf
\[
  f(t)=\mathbb P(T=t),
\]
也就是命中 target 的时间分布。主线里你看到的 \texttt{peak1}、\texttt{valley}、\texttt{peak2}，都是在这条总曲线上定义的窗口，而不是额外另造出来的对象。

\paragraph{为什么第一峰会早。}
如果一条轨迹大部分时间都待在 corridor 里，几乎不做 outside detour，也不发生明显 rollback，那么它会以较短时间到达 target。这批轨迹集中堆在较早的时间上，于是形成 \texttt{peak1}。这就是主线里说的 corridor fast branch。

\paragraph{为什么第二峰会晚。}
如果一条轨迹在 hitting 之前跑到 outside，或者更典型地，它先离开起始盆地 $A$，后来又回到 $A$，那它需要花更长时间才命中 target。这些更慢的路径质量堆在较晚时间上，于是形成 \texttt{peak2}。在当前主线语言里，第二峰的主体是 left-open detour/rollback branch，而不是“纯穿膜峰”。

\paragraph{膜穿越在这里扮演什么角色。}
膜穿越当然重要，因为它会改变进入 outside 和从 outside 回到 corridor 的难易，从而改变慢支路的组成比例。但主线现在的判断很明确：膜穿越是 slow branch 的次级修饰，而不是主导第二峰存在与否的唯一机制。

\begin{figure}[H]
  \centering
  \includegraphics[width=0.90\linewidth]{\GatingFigDir/one_target_rollback_window_bars.pdf}
  \caption{主线报告中的 rollback 窗口分解图。它直观展示了：\texttt{peak1} 主要由快支路支撑，而 \texttt{valley} 与 \texttt{peak2} 的质量更多地流向了包含 detour / rollback 的类别。}
  \label{fig:rollback-bars}
\end{figure}

\paragraph{为什么 Exact Recursion 非常适合这个例子。}
因为你不只想知道总曲线是不是双峰，还想知道：
\begin{itemize}
  \item 每个窗口的质量由哪些路径类别组成；
  \item 什么时候第一次离开 corridor；
  \item 什么时候第一次穿膜；
  \item 在 valley 和 \texttt{peak2} 里，outside 停留时间预算有多大差别。
\end{itemize}
这些都可以建立在同一个“精确推进未吸收质量”的框架上，只是多记录一些事件标签、多做一些窗口条件化统计而已。主线后续的 \path{grid2d_one_target_window_measures}、\path{grid2d_one_target_exit_timing} 和 \path{grid2d_one_target_valley_peak_budget} 都是在这条基础线上往外长出来的。

\section{AW 方法是什么，为什么很多人会把它和 Exact Recursion 放在一起比较}
\paragraph{先说最直观的区别。}
Exact Recursion 是在时间域里一轮一轮推进；AW 方法则是先到变换域，把整条时间分布的生成函数写出来，再把系数反演回时间域。所以两者都能得到 $f(t)$，但走的路径不同：
\begin{itemize}
  \item Exact Recursion：先有每一步的概率流，再拼出整条曲线；
  \item AW：先有整条曲线的生成函数，再用数值反演取出每个时间系数。
\end{itemize}

\paragraph{什么是生成函数。}
对首达 pmf
\[
  f(1),f(2),f(3),\dots,
\]
它的生成函数定义为
\begin{equation}
  \widetilde f(z)=\sum_{t\ge 1} f(t)z^t.
  \label{eq:genfun-def}
\end{equation}
你可以把它想象成“把整个时间序列压缩进一个关于 $z$ 的函数里”。一旦这个函数能被高效计算，就有机会反过来把每个系数 $f(t)$ 再取出来。

\paragraph{propagator 是什么。}
propagator 可以理解成“从一个起点到另一个位置，在给定时间或变换域下传播过去的权重”。如果在时间域里写，它长得像
\[
  P_{ij}(t)=\mathbb P_i(X_t=j).
\]
如果写成生成函数，就会变成
\[
  \widetilde P_{ij}(z)=\sum_{t\ge 0}\mathbb P_i(X_t=j)z^t.
\]
对于很多具有高对称性的模型，例如某些 ring 或 defect 很少的结构，$\widetilde P_{ij}(z)$ 可以写得很漂亮，这正是 Giuggioli/Sarvaharman 这一类 propagator-renewal 路线常常能发挥威力的原因。

\paragraph{single-target renewal 关系。}
如果目标只有一个点 $m$，那么最常见的关系是
\begin{equation}
  \widetilde f_{i\to m}(z)=\frac{\widetilde P_{im}(z)}{\widetilde P_{mm}(z)}.
  \label{eq:renewal-single}
\end{equation}
这条式子的直觉是：从 $i$ 传播到 $m$ 的所有路径，可以按“第一次命中 $m$ 是在什么时候”切开。第一次命中之后，后面在 $m$ 附近的传播由 $\widetilde P_{mm}(z)$ 接上。把这个 renewal 关系解出来，就得到单目标首达的生成函数。

\paragraph{AW/Cauchy-FFT 反演。}
一旦你有了 $\widetilde f(z)$，接下来要做的是“取出 $z^t$ 的系数”。最一般的思路是 Cauchy 系数公式：
\begin{equation}
  f(t)=[z^t]\widetilde f(z)
  =
  \frac{1}{2\pi i}\oint_{|z|=r}\frac{\widetilde f(z)}{z^{t+1}}\,dz.
  \label{eq:cauchy-coeff}
\end{equation}
AW 或离散 Cauchy-FFT 的做法，就是把这条圆周积分离散化成一圈采样点：
\begin{equation}
  f(t)
  \approx
  \frac{r^{-t}}{m}
  \sum_{k=0}^{m-1}
  \widetilde f\!\left(re^{2\pi i k/m}\right)
  e^{-2\pi i kt/m}.
  \label{eq:aw-basic}
\end{equation}
它看起来像 Fourier 系列，是因为本质上就是在圆周上做离散频率采样，再把系数倒回来。

\paragraph{这些参数是什么意思。}
\begin{itemize}
  \item $r$ 是取样圆的半径。太接近 1 时，系数提取更稳定于长尾，但数值条件可能变差；太小又会放大截断误差。
  \item $m$ 是圆周上的采样点数。采样点更多，通常反演更精细，但每次都要评估 $\widetilde f(z_k)$，总成本会上去。
  \item 如果每个 $z_k$ 上的 $\widetilde f(z_k)$ 本身要通过解一个线性系统才能得到，那真正昂贵的往往不是 FFT，而是每个采样点上的前置求解。
\end{itemize}

\paragraph{为什么 AW 经常拿来和 Exact 对比。}
因为它们解决的是同一个问题：都想得到首达时间分布 $f(t)$。所以非常自然地会问：
\begin{quote}
同一个模型下，是直接在时间域一轮一轮推更划算，还是先在变换域把 $\widetilde f(z)$ 写出来再反演更划算？
\end{quote}
你的仓库之所以要单独做 \path{luca_vs_recursion_unified_benchmark}，就是在系统地回答这个问题。

\section{Luca/GF 家族到底是什么}
\paragraph{不要把它想成一个单独公式。}
在你的仓库语境里，Luca/GF 家族更像是一整条“方法链”：
\begin{enumerate}
  \item 先写 defect-free propagator 或闭式 Green function；
  \item 再把局部 defect、door、sticky、shortcut 等有限异质性压缩成低维修正；
  \item 通过 Woodbury、determinant closure 或 resolvent 恢复 defected propagator；
  \item 再做 single-target 或 multi-target renewal，得到 first-passage 的生成函数；
  \item 最后用 AW/Cauchy-FFT 把时间系数反演出来。
\end{enumerate}

\paragraph{为什么这条路线会有时非常强。}
如果问题有很好的解析对称性，而且 defect support 很小，那么“先有一个漂亮的无缺陷传播子，再做一个低维修正”就会非常划算。因为你不需要在整个巨大状态空间上逐步推进，只需要在少量 defect 自由度上解问题，然后再把答案拼回来。

\paragraph{为什么它有时又会很贵。}
一旦 defect 不再稀疏，或者需要对很多 $z_k$ 做同样昂贵的线性求解，这条路线的成本就会迅速抬高。你的 unified benchmark 已经把这一点总结得很清楚：在有解析 PGF 的单目标 ring，或 ultra-sparse defect 的某些 two-target 工作负载里，Luca/GF 可以赢；但在二维主线、encounter、reflecting 中等或高 defect 的案例里，time recursion 往往更实用。

\begin{figure}[H]
  \centering
  \includegraphics[width=0.88\linewidth]{\BenchFigDir/unified_runtime_config_overview.pdf}
  \caption{统一 benchmark 的配置概览图。它提醒我们：仓库里并不是只研究一种模型，因此不会存在一条在所有工作负载上都最优的方法路线。}
  \label{fig:benchmark-overview}
\end{figure}

\paragraph{为什么 ring 单目标是 Luca/GF 的正锚点。}
因为那类模型往往能把生成函数写得很干净。例如 \path{ring_lazy_flux} 这条线就明确是“Chebyshev 解析生成函数 + AW 反演；flux 交叉验证”的结构。这里 AW/Luca 并不是陪跑，而是很自然的主力方法之一。

\section{Exact Recursion、AW、Luca/GF、MC 到底怎么比较}
\paragraph{先给一张总表。}
\small
\setlength{\LTleft}{0pt}
\setlength{\LTright}{0pt}
\begin{longtable}{>{\raggedright\arraybackslash}p{0.14\linewidth} >{\raggedright\arraybackslash}p{0.14\linewidth} >{\raggedright\arraybackslash}p{0.14\linewidth} >{\raggedright\arraybackslash}p{0.18\linewidth} >{\raggedright\arraybackslash}p{0.28\linewidth}}
\toprule
方法 & 直接求的对象 & 是否精确 & 主要误差来源 & 在本仓库里的典型角色 \\
\midrule
\endhead
Exact Recursion & 未吸收分布 $u_t$ 与首达 pmf $f(t)$ & 是，在给定时间窗内精确 & 浮点误差、时间截断 & 当前二维主线的默认生产器；也是 encounter 与多数 reflecting / two-target 工作负载的实务首选 \\
AW 反演 & 首达生成函数 $\widetilde f(z)$ 的时间系数 & 反演本身是数值近似，但可很高精度 & 采样点数 $m$、半径 $r$、条件数、截断 & 在有良好生成函数表达时很强；常用于和 exact 做交叉验证 \\
Luca/GF 家族 & propagator / Green function，经 renewal 得到 $\widetilde f(z)$ & 若前置解析和数值求解都正确，可高精度 & defect-reduced 解的代价、每个 $z_k$ 的线性求解、反演误差 & 单目标 ring、极稀疏 defect 的正锚点；统一 benchmark 的重要对照路线 \\
Monte Carlo & 轨迹样本的经验频率 & 否，是统计估计 & 样本噪声、尾部方差、binning/smoothing & 主要用于解释、可视化和 sanity check，而不是主曲线生产器 \\
\bottomrule
\end{longtable}
\normalsize

\paragraph{为什么你的二维主线主要用 Exact Recursion。}
因为它同时满足三件事：
\begin{enumerate}
  \item 目标是完整首达曲线和窗口结构，而不是只要一个平均值；
  \item defect / membrane / local bias 的结构已经复杂到让“每个 $z_k$ 做一轮大求解”的代价变得不划算；
  \item 你后续还要做大量 window-conditioned observables 和机制分解，而 time recursion 很自然就能和这些统计挂起来。
\end{enumerate}

\paragraph{为什么 ring 的某些解析单目标例子更适合 AW/Luca。}
因为那一类模型的闭式 propagator 和单目标 renewal 关系都很干净。一旦生成函数已经在手，反演整条曲线会非常自然。这也是为什么 unified benchmark 里 RING-1T-paper 是 Luca/GF 的正锚点。

\paragraph{为什么 MC 在你的工作里不是主生产器。}
这并不是说 MC 没用，而是它的角色不同。你的主结果经常需要：
\begin{itemize}
  \item 精确辨认双峰是否存在；
  \item 比较 valley 和 second peak 的细微差异；
  \item 稳定输出 timing ratio、outside budget、post-crossing budget；
  \item 在大量参数点上重复跑并保持一致口径。
\end{itemize}
这些任务如果全靠 MC，样本量会很大、尾部会比较吵、窗口边界也不稳定。Exact Recursion 则更适合作为 deterministic baseline，MC 更像辅助解释层。

\section{这套方法怎样贯穿你的研究主线}
\paragraph{不是只算一条曲线。}
对你的研究来说，Exact Recursion 的价值不在于“把首达曲线画出来”，而在于它提供了一块可重复使用的基底。只要你能把事件条件、区域标签或扩展状态编码进状态空间或后处理，就能在同一框架里继续计算很多机制量。

\paragraph{它怎样支撑 one-target companions。}
\begin{itemize}
  \item 在 \path{grid2d_one_target_window_measures} 里，它支撑了窗口条件化的 occupancy share 和 ever-visit probability；
  \item 在 \path{grid2d_one_target_exit_timing} 里，它支撑了 $\tau_{\mathrm{out}}$ 和 $\tau_{\mathrm{mem}}$ 的窗口条件 CDF、timing ratio 和 early/late/no-exit 分解；
  \item 在 \path{grid2d_one_target_valley_peak_budget} 里，它继续支撑 valley 与 \texttt{peak2} 的 outside-time share、outside budget 和 post-crossing budget 对比。
\end{itemize}
这些看起来是不同问题，但底层共同点都是：先得到足够精确的首达结构，再在此基础上做条件化和记账。

\paragraph{它怎样和 benchmark 连起来。}
\path{luca_vs_recursion_unified_benchmark} 做的不是推翻这条主线，而是把它放进更大的方法比较语境里：什么时候应继续用 time recursion，什么时候应该切换到 Luca/GF。也就是说，benchmark 不是反对 Exact Recursion，而是在帮你说明它的优势边界和局限边界。

\section{常见误解}
\paragraph{误解 1：你主要是不是在做 Monte Carlo？}
不是。MC 在你的仓库里主要承担辅助角色：做代表轨迹、画 overlay、做少量 sanity check。真正的大量主曲线和机制量，尤其是二维主线，并不是靠 MC 样本频率去生产的。

\paragraph{误解 2：你是不是主要用 Luca/GF？}
也不是。Luca/GF 是仓库里非常重要的一条方法线，而且在某些单目标 ring 或稀疏 defect 情况下很强，但它不是当前二维主线的默认生产器。你现在的主线，更准确地说是“以 Exact Recursion 为默认求解器，再用 AW/Luca 做局部正锚点、交叉验证和方法比较”。

\paragraph{误解 3：你的研究是不是纯 lazy？}
不是全都 lazy。仓库整体上既有 lazy 也有 non-lazy 线；只是当前二维主线常用 baseline 的确是 lazy 参数设置，例如很多例子固定 $q=0.8$。方法本身并不依赖“必须 lazy”，它只依赖你把一步转移概率写清楚。

\paragraph{误解 4：`exact' 到底精确在哪里？}
它的“精确”不是说“一定有闭式公式”，而是说：在给定有限状态空间和给定时间窗内，它直接精确推进全部概率质量，没有 Monte Carlo 采样误差。实际数值上的误差来源只剩下浮点和你主动设定的时间截断。

\section{阅读指南与最小复现入口}
\paragraph{如果你只想最快读懂主线。}
建议顺序是：
\begin{enumerate}
  \item 先读本文第 2--5 节，真正理解 Exact Recursion 在算什么；
  \item 再对照 \path{grid2d_one_two_target_gating} 的 one-target shared symmetric baseline；
  \item 然后读 \path{grid2d_one_target_exit_timing} 和 \path{grid2d_one_target_valley_peak_budget}，看同一套递推怎样延展出 timing 和 budget 量；
  \item 最后再回头看 \path{luca_vs_recursion_unified_benchmark}，理解为什么同一个仓库里需要两条方法线。
\end{enumerate}

\paragraph{与仓库代码最相关的入口。}
\begin{itemize}
  \item 二维主线状态与转移构造：\path{packages/vkcore/src/vkcore/grid2d/model_core_reflecting.py}
  \item one-target 精确递推骨架：\path{packages/vkcore/src/vkcore/grid2d/rect_bimodality/cli.py}
  \item 当前主线报告：\path{research/reports/grid2d_one_two_target_gating/}
  \item 统一方法比较：\path{research/reports/luca_vs_recursion_unified_benchmark/}
\end{itemize}

\paragraph{最小复现命令。}
对本文自身：
\begin{quote}
\path{python3 scripts/reportctl.py build --report exact_recursion_method_guide --lang cn}
\end{quote}
对主线 one-target 机制稿：
\begin{quote}
\path{python3 scripts/reportctl.py build --report grid2d_one_two_target_gating --lang cn}
\end{quote}
对统一 benchmark：
\begin{quote}
\path{python3 scripts/reportctl.py build --report luca_vs_recursion_unified_benchmark --lang cn}
\end{quote}

\begin{thebibliography}{9}
\bibitem{aw1992}
J. Abate and W. Whitt,
\emph{The Fourier-series method for inverting transforms of probability distributions},
Queueing Systems \textbf{10}, 5--88 (1992).

\bibitem{sarvaharman2020pre}
S. Sarvaharman and L. Giuggioli,
\emph{First-passage-time statistics of a random walker under biased quenched disorder},
Phys. Rev. E \textbf{102}, 062124 (2020).

\bibitem{giuggioli2022transmission}
L. Giuggioli and S. Sarvaharman,
\emph{Spatio-temporal dynamics of random transmission events: from information sharing to epidemic spread},
J. Phys. A: Math. Theor. \textbf{55}, 375005 (2022).

\bibitem{giuggioli2024multi}
L. Giuggioli, S. Sarvaharman, D. Das, D. Marris, and T. Kay,
\emph{Multi-target search in bounded and heterogeneous environments: a lattice random walk perspective},
arXiv:2311.00464v2 [cond-mat.stat-mech] (2024).
\end{thebibliography}

\end{document}
